\subsection*{一步一步建立直观}

\paragraph{1. 把系数当成两个运动旋钮}

先固定两根方向箭头，只想象改变它们前面的系数。每一个系数组合产生一个终点，所有可能终点的集合才叫张成空间。右图显示的是其中一个组合；有限次拖动只能提供例子，不能代替对所有实系数的定义。



\paragraph{2. 可到达与唯一到达是两件事}

若两个方向不共线，可以沿它们分别补足所需的两份位移。对于平面上的目标，联立两条坐标方程便得到系数；非零行列式保证答案唯一。基既要能到达全部目标，又要避免用不同配方描述同一目标。



\paragraph{3. 用退化情形检验条件}

把 $v$ 设成 $2u$，同一个终点可以通过增加一份 $u$、减少半份 $v$ 保持不变。这个看得见的冗余就是线性相关。此时系数空间仍是二维，但输出只是一条线；区分输入参数数目与实际可达维数十分重要。


